\section{Discrete-Time Hedging Problem}
\label{sec:problem}
% ---------------------------------------------------------------------------

\subsection{Trading dates, payoff, and self-financing gain}

Let
\[
  0=t_0<t_1<\cdots<t_n=T
\]
be the trading grid, and let $S_k:=S_{t_k}$ be the underlying price. The
liability is a short European call
\[
  H(S_n)=(S_n-K)_+.
\]
At $t_k$, for $k=0,\ldots,n-1$, the hedger holds $\delta_k$ units of the
underlying over $[t_k,t_{k+1}]$. Set $\delta_{-1}=0$. The self-financing gain
from the risky asset is
\[
  G_T(\delta)=\sum_{k=0}^{n-1}\delta_k(S_{k+1}-S_k).
\]
Interest rates are set to zero in all simulations.

Define the trade at date $t_k$ by
\[
  u_k:=\delta_k-\delta_{k-1},\qquad k=0,\ldots,n-1.
\]
For a general formulation, set $\delta_n=0$ and define the terminal liquidation
trade $u_n=-\delta_{n-1}$. The total trading cost is
\[
  C_T(\delta)=\sum_{k=0}^{n}c_k(u_k).
\]
Following the numerical convention used in the experiments reported below,
terminal liquidation is free: $c_n\equiv0$. Thus the risky position is valued
through its final gain over $[t_{n-1},T]$, and no additional fee is charged for
closing it at maturity. This convention must be distinguished from a model
with a positive terminal liquidation cost.

Given initial cash $p_0$, terminal profit and loss is
\begin{equation}
  \Pi_T(p_0,\delta)
  =p_0-H(S_n)+G_T(\delta)-C_T(\delta),
  \label{eq:pnl}
\end{equation}
and terminal loss is
\begin{equation}
  L_T(p_0,\delta)
  :=-\Pi_T(p_0,\delta)
  =H(S_n)-p_0-G_T(\delta)+C_T(\delta).
  \label{eq:loss}
\end{equation}
The numerical tables use $p_0=0$. They therefore report an uncentered liability
loss, not the net hedging error after receipt of an option premium.

\subsection{Transaction-cost functions}

For $k<n$, the two specifications are
\begin{align}
  c_k^{\mathrm{lin}}(u)&=\kappa S_k|u|,
  \label{eq:lin_cost}\\
  c_k^{\mathrm{quad}}(u)&=\kappa S_k u^2.
  \label{eq:quad_cost}
\end{align}
The parameter $\kappa$ is dimensionless under the chosen convention. For
proportional costs, $\kappa=0.01$ means a one-way cost equal to one percent of
the traded notional, which is high for a liquid equity and should be interpreted
as a stress level rather than as a universal market calibration.

\subsection{Translation of the loss origin}

\begin{proposition}[Cash translation]
For any deterministic cash amount $a$ and any policy class $\mathcal{D}$,
\[
  \arg\min_{\delta\in\mathcal{D}}
  \CVaR_{\alpha}\!\left(L_T(p_0+a,\delta)\right)
  =
  \arg\min_{\delta\in\mathcal{D}}
  \CVaR_{\alpha}\!\left(L_T(p_0,\delta)\right).
\]
Moreover, absolute CVaR differences between two policies are invariant to the
cash shift, whereas ratios and percentage improvements generally are not.
\end{proposition}

\begin{proof}
Equation~\eqref{eq:loss} gives
$L_T(p_0+a,\delta)=L_T(p_0,\delta)-a$. CVaR is cash-equivariant, so
$\CVaR_{\alpha}(L-a)=\CVaR_{\alpha}(L)-a$. The same constant is subtracted
from the objective for every policy and from both members of an absolute
comparison. A ratio changes because its denominator is shifted.
\end{proof}

Accordingly, the primary comparisons below are stated in absolute loss units.
Relative percentages are retained only as descriptive summaries under the
explicit convention $p_0=0$.

% ---------------------------------------------------------------------------

\section{Risk Criterion and Empirical Optimization}
\label{sec:cvar}
% ---------------------------------------------------------------------------

\subsection{Population CVaR}

For an integrable loss $L$ and confidence level $\alpha\in(0,1)$, the
Rockafellar--Uryasev representation is
\begin{equation}
  \CVaR_{\alpha}(L)
  =\min_{w\in\R}
  \left\{
    w+\frac{1}{1-\alpha}\E\!\left[\pospart{L-w}\right]
  \right\}.
  \label{eq:ru}
\end{equation}
For a continuous loss distribution, a minimizer is the usual
$\alpha$-quantile. With atoms, the minimizer is generally a quantile set, and
CVaR should not be reduced uncritically to a conditional expectation given
$L\geq\VaR_{\alpha}(L)$.

The function in Equation~\eqref{eq:ru} is convex in $w$ for a fixed policy. It
is not generally convex in the neural-network parameters.

\subsection{Finite-batch empirical CVaR}

For a batch of $B$ simulated losses $L_1(\theta),\ldots,L_B(\theta)$, define
\begin{equation}
  \widehat F_B(\theta,w)
  =w+\frac{1}{(1-\alpha)B}
  \sum_{i=1}^{B}\pospart{L_i(\theta)-w}.
  \label{eq:emp_ru}
\end{equation}
Profiling out $w$ minimizes the \emph{empirical batch objective} exactly. It
does not make the resulting random mini-batch criterion identical to the
population objective, because in general
\[
  \E_{\mathcal B}\!\left[\min_w\widehat F_{\mathcal B}(\theta,w)\right]
  \neq
  \min_w\E_{\mathcal B}\!\left[\widehat F_{\mathcal B}(\theta,w)\right].
\]
The profiled estimator is nevertheless a standard consistent empirical
approximation as the batch size increases under the usual regularity
conditions.

For an exact finite-sample implementation, sort the batch losses in descending
order,
\[
  L_{[1]}\geq L_{[2]}\geq\cdots\geq L_{[B]},
\]
and let
\[
  q=(1-\alpha)B,\qquad r=\lfloor q\rfloor,
  \qquad \lambda=q-r.
\]
Then, for $q>0$, empirical expected shortfall is
\begin{equation}
  \widehat{\CVaR}_{\alpha,B}
  =\frac{1}{q}
  \left(
    \sum_{j=1}^{r}L_{[j]}+\lambda L_{[r+1]}
  \right),
  \label{eq:exact_emp_cvar}
\end{equation}
with the fractional term omitted when $\lambda=0$. Away from ties, an exact
gradient is
\begin{equation}
  \nabla_{\theta}\widehat{\CVaR}_{\alpha,B}
  =\frac{1}{q}
  \left(
    \sum_{j=1}^{r}\nabla_{\theta}L_{[j]}
    +\lambda\nabla_{\theta}L_{[r+1]}
  \right).
  \label{eq:exact_emp_cvar_grad}
\end{equation}
At ties, this expression is interpreted as a valid subgradient after a
specified tie-breaking rule. A hard indicator that includes only losses
strictly above a sample quantile omits the fractional boundary observation
whenever $(1-\alpha)B$ is non-integer. For $B=4,096$ and $\alpha=0.95$, the
exact tail mass is $204.8$: 204 observations receive full weight and the next
observation receives weight $0.8$.

\subsection{Threshold profiling and training stalls}

If $w$ is updated jointly by stochastic gradient descent, its batch derivative
is
\[
  \frac{\partial\widehat F_B}{\partial w}
  =1-\frac{1}{(1-\alpha)B}
  \sum_{i=1}^{B}\1\{L_i>w\},
\]
away from ties. When $w$ exceeds every loss in a batch, the policy receives no
gradient through the positive-part terms, while $w$ moves only by its optimizer
step. Profiling $w$ within each batch avoids this particular zero-mask event.
This is an implementation advantage, not a proof that the stochastic
optimization problem has become convex or noiseless.


% ---------------------------------------------------------------------------

\section{Policy Class and Analytic Gradients}
\label{sec:methodology}
% ---------------------------------------------------------------------------

\subsection{State-only shared policy}

The hedge at each date is generated by a shared-weight feedforward network,
\begin{equation}
  \delta_k=g_{\theta}(x_k)
  =\delta_{\max}\tanh\!\left(f_{\theta}(x_k)\right),
  \qquad \delta_{\max}=1.5.
  \label{eq:policy}
\end{equation}
The network $f_{\theta}$ has two hidden layers with 16 $\tanh$ units per
layer. Under GBM,
\[
  x_k=(\tau_k,\log(S_k/K)),\qquad \tau_k=T-t_k,
\]
and under Heston,
\[
  x_k=(\tau_k,\log(S_k/K),v_k^+).
\]

This policy class is not recurrent. In particular, it cannot map the same
market state to different actions for different current inventories. The cost
terms couple the outputs during training, but the gradient is not a state
variable available at inference. Thus the architecture can fit a smoother or
more conservative state map, but it cannot represent a genuine inventory-
dependent no-trade region of the form
$\delta_k=\varphi(x_k,\delta_{k-1})$. This restriction is central when
interpreting the turnover results.

\subsection{Network backward pass}

For one observation, the forward pass is
\begin{align*}
  z_1&=xW_1+b_1, & a_1&=\tanh z_1,\\
  z_2&=a_1W_2+b_2, & a_2&=\tanh z_2,\\
  z_3&=a_2W_3+b_3, & \delta&=\delta_{\max}\tanh z_3.
\end{align*}
Given an upstream derivative $g_{\delta}=\partial\mathcal L/\partial\delta$,
\begin{align*}
  g_{z_3}&=g_{\delta}\,\delta_{\max}(1-\tanh^2z_3),\\
  \nabla_{W_3}\mathcal L&=a_2^{\top}g_{z_3},
  &\nabla_{b_3}\mathcal L&=g_{z_3},\\
  g_{a_2}&=g_{z_3}W_3^{\top},
  &g_{z_2}&=g_{a_2}\odot(1-a_2^2),\\
  \nabla_{W_2}\mathcal L&=a_1^{\top}g_{z_2},
  &\nabla_{b_2}\mathcal L&=g_{z_2},\\
  g_{a_1}&=g_{z_2}W_2^{\top},
  &g_{z_1}&=g_{a_1}\odot(1-a_1^2),\\
  \nabla_{W_1}\mathcal L&=x^{\top}g_{z_1},
  &\nabla_{b_1}\mathcal L&=g_{z_1}.
\end{align*}
Because the same parameters are used at every date, the parameter gradients are
summed over $k=0,\ldots,n-1$ before each optimizer update.

\subsection{Indexed transaction-cost gradient}

Let $u_j=\delta_j-\delta_{j-1}$ for $j<n$ and
$u_n=-\delta_{n-1}$. From Equation~\eqref{eq:loss},
\begin{equation}
  \frac{\partial L_T}{\partial\delta_k}
  =-(S_{k+1}-S_k)+c_k'(u_k)-c_{k+1}'(u_{k+1}),
  \qquad k=0,\ldots,n-1,
  \label{eq:delta_gradient_general}
\end{equation}
where $c_n'\equiv0$ under the free-terminal-liquidation convention. The
important index is $k+1$ in the second cost derivative. For an interior date,
the subsequent transaction cost is scaled by $S_{k+1}$, not by $S_k$.

For proportional costs and $u\neq0$,
\begin{equation}
  c_j'(u)=\kappa S_j\operatorname{sign}(u).
  \label{eq:linear_subgrad}
\end{equation}
At $u=0$, the subdifferential is
$\kappa S_j[-1,1]$; the implementation uses zero as a valid subgradient. For
quadratic costs,
\begin{equation}
  c_j'(u)=2\kappa S_j u.
  \label{eq:quadratic_grad}
\end{equation}

\subsection{Gradient checks}

Table~\ref{tab:gradcheck} reproduces the archived centered finite-difference
check for a small network using step size $10^{-5}$. The reported relative-error
convention is
\[
  \frac{|g_{\mathrm{ana}}-g_{\mathrm{FD}}|}
  {\max\{1,|g_{\mathrm{ana}}|,|g_{\mathrm{FD}}|\}}.
\]

\begin{table}[htbp]
\centering
\caption{Analytic network backward pass versus centered finite differences
($\epsilon=10^{-5}$).}
\label{tab:gradcheck}
\begin{tabular}{lcc}
\toprule
Parameter & Maximum absolute difference & Maximum reported relative error \\
\midrule
$W_1$ & $1.65\times10^{-9}$ & $1.39\times10^{-9}$ \\
$b_1$ & $3.14\times10^{-10}$ & $5.91\times10^{-11}$ \\
$W_2$ & $5.07\times10^{-11}$ & $9.35\times10^{-10}$ \\
$b_2$ & $1.30\times10^{-10}$ & $3.09\times10^{-11}$ \\
$W_3$ & $5.85\times10^{-11}$ & $1.89\times10^{-11}$ \\
$b_3$ & $7.73\times10^{-11}$ & $1.28\times10^{-11}$ \\
\bottomrule
\end{tabular}
\end{table}

The existing quadratic-cost pathwise check reports a maximum absolute deviation
of $2.5\times10^{-10}$ and a maximum relative deviation of
$1.0\times10^{-9}$. The current repository additionally tests the $S_{k+1}$ index, the
fractional empirical-CVaR boundary weight, and the inventory-aware backward
pass. An independent automatic-differentiation comparison of the complete
profiled empirical-CVaR parameter gradient remains outstanding.

\subsection{Optimization}

Parameters are updated with Adam \citep{kingma2015}, using learning rate
$3\times10^{-3}$, $\beta_1=0.9$, and $\beta_2=0.999$. Fresh Monte Carlo
batches of 4,096 paths are generated at every iteration. The reported
experiments use 1,200--1,500 updates depending on the configuration. This is
on-the-fly Monte Carlo stochastic optimization rather than optimization over a
fixed finite training set.

For exact reproducibility, the Adam stabilizer $\varepsilon$, parameter dtype,
random-number generator, seed, feature scaling, initialization distribution,
iteration count by experiment, and checkpoint-selection rule must also be
archived. These items cannot be inferred from the reported point estimates.

% ---------------------------------------------------------------------------
